\documentclass[conference]{IEEEtran}
\usepackage{amsmath,amssymb,amsfonts}
\usepackage{algorithmic}
\usepackage{algorithm}
\usepackage{graphicx}
\usepackage{textcomp}
\usepackage{xcolor}
\usepackage{booktabs}
\usepackage{multirow}
\usepackage{hyperref}

\begin{document}

\title{AdaPrune: Online Adaptive Pruning for Gradient Boosting Decision Trees}

\author{
\IEEEauthorblockN{Your Name}
\IEEEauthorblockA{Your Affiliation \\
Email: your.email@example. com}
}

\maketitle

\begin{abstract}
Gradient Boosting Decision Trees (GBDT) are powerful machine learning models widely used in various applications. However, selecting appropriate pruning parameters (e.g., max\_depth, min\_samples\_leaf) typically requires extensive hyperparameter tuning, and fixed parameters throughout training may lead to suboptimal generalization. In this paper, we propose \textbf{AdaPrune}, an online adaptive pruning framework that dynamically adjusts tree complexity during the boosting process. AdaPrune monitors the training dynamics, including the generalization gap between training and validation performance, and adaptively modifies pruning parameters to prevent overfitting while maintaining model expressiveness. We introduce multiple adaptation strategies:  Gap-based, Trend-based, Data-aware, and a Hybrid approach that combines their strengths.  Extensive experiments on 15 benchmark datasets demonstrate that AdaPrune effectively reduces overfitting (10\% lower generalization gap) while maintaining competitive predictive accuracy compared to standard GBDT implementations. 
\end{abstract}

\begin{IEEEkeywords}
Gradient Boosting, Decision Trees, Adaptive Pruning, Overfitting, Online Learning
\end{IEEEkeywords}

%==============================================================================
\section{Introduction}
%==============================================================================

Gradient Boosting Decision Trees (GBDT) \cite{friedman2001greedy} have become one of the most successful machine learning algorithms, achieving state-of-the-art performance in various domains including recommendation systems, fraud detection, and scientific applications. Popular implementations such as XGBoost \cite{chen2016xgboost}, LightGBM \cite{ke2017lightgbm}, and CatBoost \cite{prokhorenkova2018catboost} have further improved the efficiency and scalability of GBDT. 

Despite their success, GBDT models are prone to overfitting, especially when dealing with noisy data or limited training samples. The key hyperparameters controlling model complexity—such as maximum tree depth (\texttt{max\_depth}), minimum samples per leaf (\texttt{min\_samples\_leaf}), and regularization parameters—significantly impact generalization performance.  Traditional approaches rely on grid search or Bayesian optimization to find optimal parameters, which is computationally expensive and assumes that a single set of parameters is optimal throughout the entire training process.

\textbf{Key Insight:} The optimal tree complexity may vary during different stages of the boosting process.  Early iterations may benefit from more complex trees to capture the main patterns, while later iterations might require simpler trees to avoid fitting noise. 

In this paper, we propose \textbf{AdaPrune}, an online adaptive pruning framework that dynamically adjusts pruning parameters during training based on real-time monitoring of training dynamics. Our main contributions are:

\begin{enumerate}
    \item We propose an online adaptive pruning mechanism that monitors the generalization gap and adjusts tree complexity during the boosting process. 
    \item We introduce four adaptation strategies: Gap-based, Trend-based, Data-aware, and Hybrid, each suitable for different scenarios. 
    \item We provide a comprehensive data profiling module that analyzes dataset characteristics to initialize appropriate parameters.
    \item We conduct extensive experiments on 15 benchmark datasets, demonstrating that AdaPrune reduces overfitting by 10\% while maintaining competitive accuracy.
\end{enumerate}

%==============================================================================
\section{Related Work}
%==============================================================================

\subsection{Gradient Boosting Decision Trees}

Gradient Boosting \cite{friedman2001greedy} builds an ensemble of weak learners (typically decision trees) in a stage-wise fashion, where each new tree corrects the residual errors of the previous ensemble. The objective function at iteration $t$ is: 

\begin{equation}
\mathcal{L}^{(t)} = \sum_{i=1}^{n} l(y_i, \hat{y}_i^{(t-1)} + f_t(x_i)) + \Omega(f_t)
\end{equation}

where $l$ is the loss function, $f_t$ is the new tree, and $\Omega(f_t)$ is the regularization term.

Modern implementations like XGBoost \cite{chen2016xgboost} introduce additional regularization: 

\begin{equation}
\Omega(f) = \gamma T + \frac{1}{2}\lambda \sum_{j=1}^{T} w_j^2
\end{equation}

where $T$ is the number of leaves and $w_j$ are the leaf weights.

\subsection{Hyperparameter Optimization}

Traditional hyperparameter optimization methods include grid search, random search \cite{bergstra2012random}, and Bayesian optimization \cite{snoek2012practical}. These methods treat hyperparameters as fixed values optimized before training. Recent work on dynamic hyperparameter scheduling \cite{li2017hyperband} allows early stopping of poorly performing configurations but still assumes static parameters during individual training runs.

\subsection{Adaptive Learning}

Adaptive learning rate methods like Adam \cite{kingma2014adam} adjust learning rates during training based on gradient statistics. Similarly, curriculum learning \cite{bengio2009curriculum} proposes training on easier examples first.  Our work extends this adaptive philosophy to tree complexity in gradient boosting.

%==============================================================================
\section{Methodology}
%==============================================================================

\subsection{Framework Overview}

AdaPrune consists of three main components: 

\begin{enumerate}
    \item \textbf{Data Profile Analyzer}:  Analyzes dataset characteristics including sample size, feature dimensionality, class imbalance, and estimated noise level.
    \item \textbf{State Monitor}: Tracks training dynamics including training/validation performance, generalization gap, and performance trends.
    \item \textbf{Pruning Controller}: Adjusts pruning parameters based on the current state and data profile using one of the adaptation strategies.
\end{enumerate}

\begin{figure}[htbp]
\centerline{\includegraphics[width=0.48\textwidth]{framework.png}}
\caption{AdaPrune Framework Overview.  The system monitors training dynamics and adaptively adjusts tree pruning parameters during the boosting process. }
\label{fig:framework}
\end{figure}

\subsection{Data Profile Analysis}

Before training, we analyze the dataset to estimate its characteristics:

\textbf{Noise Level Estimation:} We use a 1-NN classifier with cross-validation to estimate the Bayes error rate, which indicates the noise level: 

\begin{equation}
\text{noise\_level} = 1 - \text{accuracy}_{1\text{-NN}}
\end{equation}

\textbf{Class Overlap Score:} We compute the Fisher discriminant ratio to measure class separability:

\begin{equation}
\text{overlap} = \frac{1}{1 + \frac{(\mu_1 - \mu_2)^2}{\sigma_1^2 + \sigma_2^2}}
\end{equation}

These characteristics guide the initialization of pruning parameters and the sensitivity of adaptation. 

\subsection{State Monitoring}

At each adaptation checkpoint (every $k$ iterations), we compute:

\begin{itemize}
    \item \textbf{Generalization Gap}: $g_t = \text{acc}_{\text{train}}^{(t)} - \text{acc}_{\text{val}}^{(t)}$
    \item \textbf{Gap Trend}: $\Delta g_t = g_t - g_{t-1}$
    \item \textbf{Overfitting Score}: $o_t = \sigma(g_t \cdot 5) \cdot 0.6 + \sigma(\Delta g_t \cdot 50) \cdot 0.4$
    \item \textbf{Plateau Score}:  Measures if validation performance has stagnated
\end{itemize}

\subsection{Adaptation Strategies}

\subsubsection{Gap-Based Strategy}

This strategy focuses on the current generalization gap: 

\begin{equation}
\text{max\_depth}_{t+1} = 
\begin{cases}
\text{max\_depth}_t - 1 & \text{if } g_t > \theta_g \\
\text{max\_depth}_t + 1 & \text{if } o_t < 0.3 \text{ and underfitting} \\
\text{max\_depth}_t & \text{otherwise}
\end{cases}
\end{equation}

where $\theta_g$ is the gap threshold (default:  0.05).

\subsubsection{Trend-Based Strategy}

This strategy monitors performance trends over a sliding window:

\begin{equation}
\text{adjust} = 
\begin{cases}
\text{increase regularization} & \text{if plateau detected} \\
\text{maintain} & \text{otherwise}
\end{cases}
\end{equation}

\subsubsection{Data-Aware Strategy}

This strategy incorporates dataset characteristics:

\begin{equation}
\text{strength} = (1 + \text{noise} \cdot 2) \cdot \sqrt{\frac{1000}{\max(n, 100)}} \cdot o_t
\end{equation}

Higher noise levels and smaller sample sizes lead to more aggressive regularization.

\subsubsection{Hybrid Strategy}

The hybrid strategy combines the above approaches based on training progress:

\begin{equation}
\text{strategy} = 
\begin{cases}
\text{Data-Aware} & \text{if progress} < 0.3 \\
\text{Gap-Based} & \text{if } 0.3 \leq \text{progress} < 0.7 \\
\text{Trend-Based} & \text{if progress} \geq 0.7
\end{cases}
\end{equation}

\subsection{Parameter Adjustment}

When overfitting is detected, we adjust multiple parameters simultaneously:

\begin{itemize}
    \item Decrease \texttt{max\_depth} by 1-2 levels
    \item Increase \texttt{min\_child\_weight} by factor of $1 + 0.3 \cdot \text{strength}$
    \item Increase $\gamma$ (minimum loss reduction) by $0.1 \cdot \text{strength}$
    \item Increase $\lambda$ (L2 regularization) by factor of $1 + 0.15 \cdot \text{strength}$
    \item Decrease \texttt{subsample} if severity is high
\end{itemize}

%==============================================================================
\section{Experiments}
%==============================================================================

\subsection{Experimental Setup}

\textbf{Datasets:} We evaluate on 15 benchmark datasets from OpenML, covering various domains and sizes (Table \ref{tab:datasets}).

\begin{table}[htbp]
\caption{Dataset Statistics}
\label{tab:datasets}
\centering
\begin{tabular}{lccc}
\toprule
Dataset & Samples & Features & Classes \\
\midrule
adult & 48,842 & 14 & 2 \\
bank-marketing & 45,211 & 16 & 2 \\
electricity & 45,312 & 8 & 2 \\
magic & 19,020 & 10 & 2 \\
eeg-eye-state & 14,980 & 14 & 2 \\
satimage & 6,430 & 36 & 6 \\
waveform & 5,000 & 40 & 3 \\
phoneme & 5,404 & 5 & 2 \\
spambase & 4,601 & 57 & 2 \\
segment & 2,310 & 19 & 7 \\
credit-g & 1,000 & 20 & 2 \\
vehicle & 846 & 18 & 4 \\
diabetes & 768 & 8 & 2 \\
ionosphere & 351 & 34 & 2 \\
sonar & 208 & 60 & 2 \\
\bottomrule
\end{tabular}
\end{table}

\textbf{Baselines:}
\begin{itemize}
    \item \textbf{XGB\_default}: XGBoost with default parameters
    \item \textbf{XGB\_tuned}: XGBoost with tuned parameters (max\_depth=6, min\_child\_weight=3)
    \item \textbf{LGBM\_default}:  LightGBM with default parameters
    \item \textbf{RF\_default}: Random Forest with default parameters
\end{itemize}

\textbf{Evaluation: } 5-fold stratified cross-validation.  Metrics:  Accuracy, F1-score, ROC-AUC, and Generalization Gap.

\subsection{Main Results}

\begin{table}[htbp]
\caption{Overall Performance Comparison}
\label{tab:main_results}
\centering
\begin{tabular}{lcccc}
\toprule
Method & Accuracy & F1 & ROC-AUC & Gap \\
\midrule
LGBM\_default & 0.8762 & 0.8264 & 0.9362 & 0.1140 \\
XGB\_default & 0.8740 & 0.8248 & 0.9357 & 0.1204 \\
RF\_default & 0.8739 & 0.8199 & \textbf{0.9382} & 0.1180 \\
XGB\_tuned & 0.8700 & 0.8188 & 0.9337 & 0.1167 \\
\midrule
AdaPrune\_hybrid & 0.8578 & 0.8033 & 0.9288 & \textbf{0.0900} \\
AdaPrune\_gap & 0.8580 & 0.8039 & 0.9269 & 0.0920 \\
AdaPrune\_trend & 0.8581 & 0.8037 & 0.9270 & 0.0910 \\
AdaPrune\_data & 0.8575 & 0.8026 & 0.9290 & 0.0915 \\
\bottomrule
\end{tabular}
\end{table}

\textbf{Key Observations:}
\begin{enumerate}
    \item AdaPrune achieves \textbf{21\% lower generalization gap} compared to XGB\_default (0.09 vs 0.12).
    \item The accuracy difference is modest (1.5-2\%), indicating AdaPrune maintains competitive performance while significantly reducing overfitting. 
    \item The Hybrid strategy provides the best balance between accuracy and overfitting control.
\end{enumerate}

\subsection{Scenario Analysis}

We further analyze performance on synthetic datasets with controlled characteristics:

\begin{table}[htbp]
\caption{Performance on Different Scenarios}
\label{tab:scenarios}
\centering
\begin{tabular}{lccc}
\toprule
Scenario & Best Method & AdaPrune Gap & XGB Gap \\
\midrule
Clean (large) & XGB\_default & 0.043 & 0.030 \\
Clean (small) & XGB\_tuned & 0.100 & 0.110 \\
Noisy (large) & XGB\_weak & \textbf{0.055} & 0.094 \\
Noisy (small) & XGB\_default & \textbf{0.112} & 0.200 \\
High-dim & XGB\_default & \textbf{0.160} & 0.180 \\
Imbalanced & XGB\_default & 0.040 & 0.035 \\
\bottomrule
\end{tabular}
\end{table}

AdaPrune shows particular advantages in: 
\begin{itemize}
    \item \textbf{Noisy data}: 40-45\% lower gap
    \item \textbf{Small samples}: Better generalization
    \item \textbf{High-dimensional data}: Reduced overfitting
\end{itemize}

\subsection{Ablation Study}

\begin{table}[htbp]
\caption{Ablation Study Results}
\label{tab:ablation}
\centering
\begin{tabular}{lcc}
\toprule
Configuration & Accuracy & Gap \\
\midrule
Full\_hybrid & \textbf{0.8073} & \textbf{0.0748} \\
DataAware\_only & 0.8056 & 0.0758 \\
Gap\_only & 0.8048 & 0.0773 \\
No\_adaptation & 0.8042 & 0.0900 \\
Trend\_only & 0.8042 & 0.0900 \\
\bottomrule
\end{tabular}
\end{table}

The ablation study confirms that: 
\begin{enumerate}
    \item Adaptive pruning reduces gap by 17\% compared to no adaptation
    \item The Hybrid strategy combining multiple approaches works best
    \item Data-aware initialization provides additional benefits
\end{enumerate}

%==============================================================================
\section{Discussion}
%==============================================================================

\subsection{When Does AdaPrune Help?}

Our experiments suggest AdaPrune is most beneficial when:
\begin{itemize}
    \item The dataset contains significant noise
    \item Training samples are limited
    \item Features are high-dimensional
    \item Default hyperparameters tend to overfit
\end{itemize}

\subsection{Computational Overhead}

AdaPrune introduces minimal overhead through:
\begin{itemize}
    \item Data profiling:  One-time cost at initialization
    \item State monitoring: O(1) per iteration
    \item Parameter adjustment: O(1) per adaptation round
\end{itemize}

The primary cost difference comes from using incremental training (which is necessary for adaptive adjustment) rather than optimized batch training in XGBoost. 

\subsection{Limitations}

\begin{enumerate}
    \item Current implementation is slower than optimized XGBoost/LightGBM
    \item Accuracy is slightly lower on clean, large datasets
    \item Hyperparameters of the adaptation mechanism itself require tuning
\end{enumerate}

%==============================================================================
\section{Conclusion}
%==============================================================================

We presented AdaPrune, an online adaptive pruning framework for Gradient Boosting Decision Trees.  By monitoring training dynamics and dynamically adjusting tree complexity, AdaPrune effectively reduces overfitting while maintaining competitive predictive performance.  Our experiments demonstrate a 21\% reduction in generalization gap compared to standard GBDT implementations. 

Future work includes: 
\begin{itemize}
    \item Integration with XGBoost/LightGBM native implementations for better efficiency
    \item Automatic tuning of adaptation hyperparameters
    \item Extension to regression tasks and other loss functions
    \item Theoretical analysis of convergence properties
\end{itemize}

\section*{Acknowledgment}
[Your acknowledgments here]

\bibliographystyle{IEEEtran}
\bibliography{references}

\end{document}